% Fallbacks and common problems of the Day 2 lab. The rows come from
% Labs/day02/README.md.
\begin{frame}{If a part does not work}
  \small
  \renewcommand{\arraystretch}{1.3}
  \begingroup\centering
  \begin{tabular}{@{}L{0.33\textwidth}L{0.60\textwidth}@{}}
    \toprule
    \textbf{Problem} & \textbf{Fallback} \\
    \midrule
    MicroPython does not start on your kit
      & Upload the Arduino sketch \code{sketches/imu\_data\_collection}. It
        sends the same lines. Continue with the rate check and with
        Part C. \\
    The board does not work at all
      & Start \code{python3 host/sim\_board.py --label lift --session s1}.
        Use \code{--udp 5005} in place of \code{--port PORT}. The signals
        are simulated. \\
    The lab Wi-Fi does not work
      & Do not do the Wi-Fi step of Part B. Write the reason in the
        report. \\
    Your group has no dataset at the end
      & The notebook makes the fallback dataset. Write in the data card
        that the signals are simulated. \\
    \bottomrule
  \end{tabular}\par\endgroup
  \vskip 0.6em
  \keypoint{A fallback keeps your group in the lab. Write in the report
    which fallback you used.}
\end{frame}

\begin{frame}{Common problems}
  \footnotesize
  \renewcommand{\arraystretch}{1.25}
  \setlength{\tabcolsep}{4pt}
  \begingroup\centering
  \begin{tabular}{@{}L{0.40\textwidth}L{0.55\textwidth}@{}}
    \toprule
    \textbf{Problem} & \textbf{Fix} \\
    \midrule
    \code{esptool} cannot connect
      & Start bootloader mode again. Hold the button before you connect the
        cable. \\
    \code{mpremote} prints \code{could not enter raw repl}
      & \code{main.py} uses the port. Open the REPL, press \code{Ctrl-C},
        then \code{Ctrl-]}. \\
    \code{ImportError: no module named 'lsm6ds3'}
      & Copy the driver to the board: \code{mpremote connect PORT fs cp
        board/lsm6ds3.py :lsm6ds3.py} \\
    \code{check\_rate.py} prints \code{Not enough samples: 0}
      & Check that \code{main.py} is on the board. Close the REPL and the
        Serial Monitor. \\
    \code{check\_rate.py --udp} gets no samples
      & Check the laptop address in \code{config.py}. Permit the UDP port
        5005 in the firewall of the laptop. \\
    The board does not connect to the Wi-Fi
      & Connect the antenna. Check \code{config.py}. The XIAO has 2.4 GHz
        Wi-Fi only. \\
    The logger prints \code{WARNING: samples are missing}
      & Record again. Close other programs that use the port. \\
    \bottomrule
  \end{tabular}\par\endgroup
  \vskip 0.4em
  \small
  The README of the lab has the complete list.
\end{frame}
